\documentclass[12pt,a4paper]{article}
\usepackage[utf8]{inputenc}
\usepackage[T1]{fontenc}
\usepackage[french]{babel}
\usepackage{amsmath,amssymb,amsthm}
\usepackage{geometry}
\usepackage{hyperref}
\geometry{margin=2.5cm}

\newtheorem{theoreme}{Théorème}
\newtheorem{proposition}{Proposition}
\newtheorem{lemme}{Lemme}
\newtheorem{corollaire}{Corollaire}
\newtheorem{axiome}{Axiome}
\newtheorem{definition}{Définition}

\title{%
  \textbf{Résumé de la Théorie Complète}\\[0.5em]
  \Large L'Univers est au Carré:\\
  Synthèse Mathématique et Validations Isabelle/HOL\\[0.5em]
  \large Philippe Thomas Savard, 2026
}
\author{Philippe Thomas Savard}
\date{2026}

\begin{document}
\maketitle
\tableofcontents
\newpage

\begin{abstract}
Ce document constitue un résumé de la théorie complète de l'univers est au carré,
développée par Philippe Thomas Savard depuis 2016. Il présente en condensé les quatre
chapitres complétés, les cinq validations formelles Isabelle/HOL, et la réponse
de l'auteur à l'hypothèse de Riemann. Ce document sert de point d'entrée pour les
lecteurs souhaitant avoir une vue d'ensemble avant d'approfondir chaque chapitre.
\end{abstract}

\section{Vue d'Ensemble de la Théorie}

La théorie mathématique de l'univers est au carré de Savard repose sur quatre chapitres
principaux:

\begin{enumerate}
  \item \textbf{Géométrie du spectre des nombres premiers} (≈55--60\% du contenu total).
  \item \textbf{Mécanique harmonique du chaos discret}.
  \item \textbf{Le postulat de l'univers est au carré — Squaring}.
  \item \textbf{Axiomatisation et conclusion — La quadrature par la pesée d'Archimède}.
\end{enumerate}

\section{Les Suites A et B: Fondements de la Théorie}

Toute la théorie repose sur deux suites fondamentales:

\[
  A = (a_i)_{i \geq 1}, \quad a_i = 2i - 1
\]
\[
  B = (b_j)_{j \geq 1}, \quad b_j \in \{6j-1, 6j+1\}
\]

Ces suites génèrent le \emph{spectre des nombres premiers}, une structure géométrique
dans $\mathbb{Z}$ qui encode tous les nombres premiers.

\section{L'Observation Centrale}

\begin{center}
\large\textbf{Observation fondamentale de Savard:}\\[0.5em]
\textit{«\, Quand $n \geq 1$ et que $n \leq -1$, tous les $n$ ramènent à un nombre premier $P$.}\\
\textit{Toutes les valeurs de $n$ sont la conséquence directe}\\
\textit{de la quantité de termes contenus dans les suites $A$ et $B$.}\\
\textit{Tous les $P$ entre eux respectent le rapport $1/k$.}\\
\textit{Il est numériquement valide, algébriquement incohérent.\,»}
\end{center}

\section{La Réponse à l'Hypothèse de Riemann}

L'hypothèse de Riemann (l'un des sept problèmes du prix du millénaire de l'Institut Clay)
affirme que tous les zéros non triviaux de la fonction zêta $\zeta(s)$ ont leur partie
réelle égale à $1/2$.

La réponse de Savard: la symétrie du spectre des nombres premiers entre les branches
$n \geq 1$ et $n \leq -1$, confirmée par la méthode de la pesée d'Archimède, impose
que le centre de masse spectral — correspondant à la droite critique $\Re(s) = 1/2$ —
est l'unique axe de symétrie des zéros de $\zeta(s)$.

\section{Les Cinq Validations Isabelle/HOL}

\begin{enumerate}
  \item \textbf{PrimeSpectrumGeometry.thy}: Géométrie spectrale et suites $A$, $B$.
  \item \textbf{RiemannZetaConjecture.thy}: Axiomatisation de la conjecture de Riemann.
  \item \textbf{HarmonicMechanics.thy}: Mécanique harmonique du chaos discret.
  \item \textbf{UniversSquaring.thy}: Postulat du squaring universel.
  \item \textbf{ArchimedesParabola.thy}: Méthode de la pesée d'Archimède.
\end{enumerate}

\section{Structure du Dépôt}

Le dépôt GitHub contient 19 fichiers certifiés (.tex, .thy, .pdf) validés par
GitHub Actions:
\begin{itemize}
  \item 7 fichiers \LaTeX{} (.tex): sources des documents en français (5) et anglais (2).
  \item 5 fichiers Isabelle/HOL (.thy): les validations formelles.
  \item 7 fichiers PDF (.pdf): les documents compilés et certifiés.
\end{itemize}

\section{Conclusion}

La théorie de l'univers est au carré représente dix ans de recherche mathématique
originale (2016--2026). Elle offre une réponse élégante et géométrique à la plus
célèbre des conjectures non résolues en mathématiques, et démontre que la nature
«carrée» de l'univers est la clé pour comprendre la distribution des nombres premiers.

\bigskip
\noindent\textbf{Auteur:} Philippe Thomas Savard\\
\noindent\textbf{Depuis:} 2016\\
\noindent\textbf{Dépôt:} \url{https://github.com/2racinede4carreunivers-dev/Theorie-mathematique-philippe-thomas-savard-2026}

\end{document}
